\documentclass[11pt]{article}
\usepackage[a4paper,margin=1in]{geometry}
\usepackage{amsmath,amssymb}
\usepackage[hidelinks]{hyperref}
\usepackage{booktabs,longtable,array}
\pdftrailerid{}

\newcommand{\Q}{\mathbb Q}

\title{Supplementary Tables for\\
\emph{Effective Archimedean Stark Theorems over Real Quadratic Fields}}
\author{Hainan Zhao}
\date{30 July 2026}

\begin{document}
\hbadness=10000
\maketitle

\section*{Scope and archive}

This supplement contains the record map and the complete
Artin-label interval replay referenced by the main paper.  It also
records queue-level Engine-A statistics that are not used in the
uniform theorem.

\paragraph{Archive DOI.}
\url{https://doi.org/10.5281/zenodo.21712478}.  The main paper PDF and
\TeX\ source, this supplement, and the frozen companion archive are
top-level files in the public record.

\section*{Supplementary Table S1: certificate record map}

\small
\begin{longtable}{@{}p{0.29\linewidth}
 >{\raggedright\arraybackslash}p{0.65\linewidth}@{}}
\toprule
result & record\\
\midrule
\endfirsthead
\toprule
result & record\\
\midrule
\endhead
first order six & \path{data/q7-p7-case-v1.json}\\
second order six & \path{data/q14-p7-case-v1.json}\\
RQ-000108 & \path{data/rq000108-case-v1.json}\\
RQ-000021 & \path{data/rq000021-case-v1.json}\\
ramified-prime-\(3\) control & \path{data/q57-norm27-case-v1.json}\\
RQ-002955 & \path{data/rq002955-case-v1.json}\\
order ten & \path{data/q33-p11-order10-case-v1.json}\\
cyclic-quartic character selection &
\path{artifacts/engine-c-character-selection-v1.json}\\
\(\Q(\sqrt{35})\) CM packet &
\path{artifacts/engine-c-w3-tranche-01-verified-v1.json}\\
\(e=6\) powered-unit record &
\path{artifacts/engine-c-e6-tranche-01-verified-v1.json}\\
\(e=6\) primitive correction &
\path{artifacts/engine-c-e6-primitive-packet-correction-v1.json}\\
\(\Q(\sqrt6)\), \(e=8\) proof route &
\path{data/q6-norm8-case-v3.json}\\
Engine-C scope correction &
\path{artifacts/engine-c-claim-scope-correction-v1.json}\\
Engine-C Fourier-sign correction &
\path{artifacts/engine-c-fourier-convention-correction-v1.json}\\
dual-route packet & \path{data/rq000458-dual-case-v1.json}\\
all archimedean places &
\path{artifacts/engine-b-archimedean-place-audit-v1.json}\\
quadratic theorem & \path{data/engine-a-uniform-theorem-v1.json}\\
quadratic Euler degeneracy audit &
\path{artifacts/engine-a-euler-degeneracy-v1.json}\\
Shintani source map &
\path{artifacts/shintani-1978-source-map-v1.json}\\
index-parity theorem &
\path{artifacts/results-paper-index-parity-lemma-v1.json}\\
parity audit &
\path{artifacts/results-paper-odd-index-parity-audit-v1.json}\\
literature perimeter & \path{data/literature-perimeter-v1.json}\\
Roblot sextic-overlap audit &
\path{artifacts/roblot-sextic-overlap-audit-v1.json}\\
\bottomrule
\end{longtable}

\section*{Supplementary Table S3: relation to Roblot's sextic theorem}

Here \(H/K\) is the cyclic sextic one-place ray field.  The columns
record Roblot's hypotheses (A1)--(A3), the certified class number of
\(H\), and applicability of \cite[Thm.~7.1]{Roblot2013}.  Every class
number was computed with \texttt{bnfinit(P,1)} and certified by
\texttt{bnfcertify(bnf)=1}; the local assertion was checked in
\(H/H^+\) by exact prime decomposition.  Roblot's (A4) is not part
of the existence statement in Theorem~7.1; as his footnote to (P2)
explains, (A4) is needed instead to deduce (P2) from a conjecturally
given Stark unit.  It is therefore neither assumed nor silently
included in the applicability column below.

\begin{center}
\begin{tabular}{@{}lccccc@{}}
\toprule
case & A1 & A2 & A3 & \(h_H\) & Thm.~7.1\\
\midrule
RQ-000190 & yes & yes & yes & \(1\) & applies\\
RQ-000419 & yes & yes & yes & \(1\) & applies\\
RQ-000021 & yes & yes & yes & \(1\) & applies\\
RQ-002057 & yes & yes & yes & \(1\) & no: wild at \(3\)\\
RQ-002955 & yes & yes & yes & \(1\) & applies\\
\bottomrule
\end{tabular}
\end{center}

For RQ-002057 the relative ramification index above \(3\) is \(6\).
For the other four cases no prime above \(3\) ramifies in \(H/K\).

\section*{Supplementary Table S2: complete Artin-label interval replay}

For each cyclic ray group below, \(r\) denotes the power of the
generator in the pinned PARI ray coordinates.  Each rational interval
contains exactly one positive root of the absolute packet polynomial.
\begin{longtable}{@{}p{0.14\linewidth}p{0.80\linewidth}@{}}
\toprule
case & \(r:\) isolating interval\\
\midrule
\endfirsthead
\toprule
case & \(r:\) isolating interval\\
\midrule
\endhead
RQ-000190 &
\(0:(5.325,5.326),\ 1:(2.713,2.714),\
2:(2.251,2.252),\ 3:(.187,.188),\
4:(.368,.369),\ 5:(.444,.445)\)\\
RQ-000419 &
\(0:(20.431,20.432),\ 1:(.189,.190),\
2:(1.090,1.091),\ 3:(.048,.049),\
4:(5.289,5.290),\ 5:(.916,.917)\)\\
RQ-000108 &
\(0:(7.892,7.893),\ 1:(2.242,2.243),\
2:(.126,.127),\ 3:(.445,.446)\)\\
RQ-000021 &
\(0:(6.109,6.110),\ 1:(7.496,7.497),\
2:(4.351,4.352),\ 3:(.163,.164),\
4:(.133,.134),\ 5:(.229,.230)\)\\
RQ-002057 &
\(0:(5.709,5.710),\ 1:(.028,.029),\
2:(.035,.036),\ 3:(.175,.176),\
4:(34.732,34.733),\ 5:(27.792,27.793)\)\\
RQ-002955 &
\(0:(18.232,18.233),\ 1:(2.180,2.181),\
2:(8.620,8.621),\ 3:(.054,.055),\
4:(.458,.459),\ 5:(.116,.117)\)\\
RQ-001107 &
\(0:(8.645,8.646),\ 1:(.276,.277),\
2:(1.245,1.246),\ 3:(.571,.572),\
4:(.236,.237),\ 5:(.115,.116),\
6:(3.611,3.612),\ 7:(.802,.803),\
8:(1.748,1.749),\ 9:(4.234,4.235)\)\\
\bottomrule
\end{longtable}

\section*{Engine-A queue statistics}

The uniform theorem does not depend on a finite search range.  In the
frozen verification queue there are 1,560 rows with nonempty
differenced character support, comprising 2,232 supported quadratic
character occurrences across 912 distinct character fields.  Exact
imprimitive Euler-factor evaluation finds 672 zero Euler products,
affecting 603 rows.  In 346 rows every supported derivative vanishes,
so the theorem's product is empty and \(X_A=1\).

\begin{thebibliography}{1}
\bibitem{Roblot2013}
X.-F.~Roblot, \emph{Index formulae for Stark units and their
solutions}, Pacific J.\ Math.\ \textbf{266} (2013), 391--422,
\url{https://doi.org/10.2140/pjm.2013.266.391}.
\end{thebibliography}

\end{document}
